\bookchapter{Lampiran E: Ucapan Terima Kasih dan Keterangan}

Terima kasih sudah tekun menyelesaikan “kelas ponsel” ini! Selama proses penulisan, kami merujuk banyak buku panduan resmi pabrikan, manual pengguna, serta artikel edukasi teknologi yang terbuka untuk umum demi memastikan keakuratan informasi.

Teknologi berkembang begitu cepat; isi buku ini didasarkan pada teknologi dan produk yang umum digunakan saat naskah ditulis. Saat kamu membacanya, mungkin saja sudah ada fitur dan produk baru yang bermunculan. Semoga cara berpikir dan metode yang diajarkan buku ini membantumu menghadapi berbagai gawai digital baru di masa depan dengan tenang.

\booksection{Poin Pembaruan Versi 1.6}

\begin{itemize}
  \item Memperluas "nilai aset digital dan prinsip pengelolaannya", menambah bab "pengelolaan sistem file dan gambar" (penamaan waktu+kejadian, pengelompokan, pembersihan otomatis 72 jam, penyamaran data EXIF).
  \item Menyempurnakan sistem perlindungan "Mode Bersih/Perlindungan" dan "Parasit Elektronik".
  \item Menambahkan pola penipuan umum di era AI (wajah tiruan dan suara palsu, penyamaran sebagai lembaga resmi, investasi fiktif) beserta strategi perlindungan untuk mengatasinya.
  \item Menulis ulang panduan perlindungan dari doxing dan serangan rekayasa sosial.
  \item Memperluas bab aksesibilitas (operasi lengkap VoiceOver/TalkBack, pembesaran dan filter warna, deskripsi gambar berbasis AI, teks langsung/real-time, saluran ganda panggilan RTT/TTY).
  \item Menulis ulang kerangka prompt dengan pendekatan "Peran-Latar Belakang-Tugas-Batasan", dilengkapi 3 contoh skenario lengkap.
  \item Menambahkan tiga lampiran teknologi tingkat lanjut: Tambahan Bab 1: Penjelasan perluasan memori ponsel Android (definisi hakiki, prinsip zRAM/swap, hambatan kecepatan, batas skenario penggunaan, penamaan tiap pabrikan dan cara mengaktifkannya).
  \item Tambahan Bab 20: Termux—lingkungan terminal Linux portabel di perangkat Android (fitur inti, ekosistem plugin, instalasi dan konfigurasi, serta skenario khas seperti pengembangan mobile, operasional server, dan pengujian keamanan).
  \item Tambahan Bab 23: Panduan Pemula GitHub—dari hosting kode hingga platform kolaborasi (membedakan Git vs GitHub, konsep inti, cara mendaftar dan memulai, alur kolaborasi Pull Request, pengingat keamanan, dan solusi akses dari dalam negeri).
\end{itemize}

\booksection{Poin Pembaruan Versi 1.7}

1. Empat bab baru

\begin{itemize}
  \item Bab 29: Kesehatan Digital—analisis waktu penggunaan layar, batas App, mode Jangan Ganggu, aturan 20-20-20 untuk kesehatan mata, penyelamatan leher, manajemen tidur, kecemasan informasi dan "kebutaan selektif", serta tantangan kesehatan digital seminggu.
  \item Bab 30: Screencasting dan Mobil—perbandingan tiga kubu AirPlay / Miracast / DLNA, solusi screencasting berkabel, mode desktop Samsung DeX, empat solusi koneksi mobil CarPlay / Android Auto / HiCar / CarLife, serta praktik screencasting di ruang rapat.
  \item Bab 31: Pembuatan Video Pendek—tiga pertanyaan pemilihan topik, 5 langkah pengaturan syuting, lima teknik gerak kamera (mendorong, menarik, mengayun, menggeser, mengikuti), alur kerja enam langkah CapCut, kunci pembuatan sampul, serta prinsip konten "3 detik pertama menentukan hidup-mati".
  \item Bab 32: Panduan P3K Keamanan—SOP 30 menit emas saat ponsel hilang, alur penanganan 72 jam setelah tertipu, pemeriksaan lewat safe mode saat terinfeksi malware, enam langkah merebut kembali akun yang dibajak, cara menghadapi pesan pemerasan, serta lima hal pencegahan harian.
\end{itemize}

2. Dua bab baru

\begin{itemize}
  \item Bab 27: Ramah Lansia dan Kebersamaan Keluarga—tabel perbandingan mode sederhana antar merek, enam modus penipuan dan tiga lini pertahanan, bantuan jarak jauh (Huawei MeeLink/Sunflower), mode Peduli WeChat dan input suara, serta SOS / deteksi jatuh / pengingat minum obat.
  \item Bab 28: Panduan Membeli Ponsel—metode memilih ponsel tiga langkah, tabel spesifikasi "cukup perhatikan lima hal", rekomendasi lima tingkat anggaran, tabel pemilihan untuk tujuh skenario, lima jebakan saat membeli, serta analisis waktu dan saluran pembelian.
\end{itemize}

3. Fitur pembacaan dan interaksi baru

\begin{itemize}
  \item Bilah kontrol pembacaan melayang: mendukung putar/jeda, bagian sebelumnya/berikutnya, penyesuaian kecepatan 0.5x-2.5x, dan berhenti.
  \item Tombol headset pintas di sisi kanan: memulai pembacaan dengan satu ketukan, otomatis mulai dari posisi baca saat ini.
  \item Pintasan keyboard: R untuk mengaktifkan/mematikan pembacaan, ESC untuk berhenti.
  \item Pengalaman pembacaan lebih baik: otomatis berlanjut ke bab berikutnya, dan kontrol pembacaan otomatis disembunyikan saat mencetak.
\end{itemize}

4. Pembaruan struktur

Daftar isi, jalur belajar (empat langkahlima langkah), peta bab, nomor versi (1.61.7), penomoran lampiran (limaenam), dan metadata semuanya diperbarui serentak.

5. Pembaruan kekinian konten

\begin{itemize}
  \item Standar prosesor: Snapdragon 8 Elite Gen5 / Dimensity 9500 masuk kelas papan atas (patokan tahun 2026).
  \item Referensi ponsel unggulan: iPhone 17 Pro, Samsung S26 Ultra, Xiaomi 17 Ultra, Huawei Pura 80 Ultra, vivo X300 Pro.
  \item Bab AI: diperbarui ke DeepSeek R1, ChatGPT-5, serta standar tenaga komputasi NPU di perangkat 80 TOPS.
  \item Layar lipat: rentang harga menyentuh kelas 4000-5000 yuan.
  \item Panduan membeli: seluruh pembagian anggaran dan rekomendasi skenario dimutakhirkan ke kondisi pasar tahun 2026.
  \item Jaringan dan ekosistem: Wi-Fi 7 / 5G-A / ekosistem HarmonyOS / LE Audio semuanya diperbarui ke spesifikasi terbaru.
\end{itemize}

6. Modul fitur baru: Dasbor Belajar (berada di bawah area Hero)

\begin{itemize}
  \item Diagram cincin progres baca: menampilkan persentase bab yang sudah dibaca secara real-time, dengan animasi cincin SVG.
  \item Seari mandiri kesehatan ponsel: 8 item daftar periksa keamanan (kata sandi, autentikasi dua faktor, temukan ponsel, cadangan cloud, cadangan WeChat, keamanan pembayaran, aplikasi anti-penipuan, dan pembersihan tanpa kata sandi).
  \item Persistensi status: status centang disimpan ke localStorage, skor kesehatan dihitung secara dinamis dan ditampilkan dengan bilah progres gradien.
  \item Pesan berjenjang: menampilkan umpan balik berbeda sesuai skor, seperti "sempurna/bagus/cukup/ada risiko keamanan".
\end{itemize}

\booksection{Poin Pembaruan Versi 1.8}

Menambahkan Bagian Ketujuh “Praktik Pemrograman AI” (Bab 37-47, total 11 bab): revolusi paradigma dan Agentic Engineering, lima prinsip berpikir pemrograman AI, context engineering dan specification-driven development (SDD), requirements engineering, panorama alat pemrograman AI 2026, panorama model besar dan praktik API, conversation engineering, alur lengkap proyek dari 0 ke 1, jaminan kualitas, jalan menuju tingkat lanjut (kolaborasi MCP/multi-Agent), serta buku panduan cepat. Daftar isi, jalur belajar (lima langkahenam langkah), peta bab, dan indeks panduan cepat diperbarui serentak. Konten bersumber dari “Metodologi Berpikir Pemrograman AI: Panduan dari Nol hingga Praktik Full Stack” edisi Agustus 2026.

Menambahkan Bagian Kedelapan “Praktik Pengembangan Android” (Bab 48-58, total 11 bab) dan Bagian Kesembilan “Praktik Pengembangan iOS” (Bab 59-69, total 11 bab): kedua platform mencakup penyiapan lingkungan (Android Studio / Xcode), kursus kilat bahasa (Kotlin / Swift), framework antarmuka (Jetpack Compose / SwiftUI), data dan jaringan, dua proyek praktik lengkap (aplikasi pembukuan + cek cuaca), pengujian dan penerbitan (Google Play / App Store), serta buku panduan cepat. Daftar isi, jalur belajar (enam langkahdelapan langkah), peta bab, dan indeks panduan cepat diperbarui serentak. Konten bersumber dari “Pengembangan Android dari Nol” dan “Pengembangan iOS dari Nol” edisi Agustus 2026.

\booksection{Tambahan Konten: Aplikasi Berguna dan Penyaringan Gangguan}

\begin{itemize}
  \item Bab 13 menambahkan bagian keenam "Perkakas Sistem yang Berguna": mengulas aplikasi pembersih seperti Qingzhuo, pembersihan penyimpanan bawaan sistem, National Anti-Fraud Center, Phone Manager bawaan merek, anti-gangguan dari operator, dan AdGuard, serta memberikan prinsip pemilihan "utamakan bawaan sistem / tambahan resmi / hemat dengan aplikasi pihak ketiga".
  \item Tabel perbandingan penyaringan gangguan bawaan per merek: mencakup jalur menuju fitur dan item yang disarankan diaktifkan untuk iOS, Xiaomi/Redmi, Huawei/HONOR, OPPO/realme/OnePlus, vivo/iQOO, dan Samsung.
  \item Perluasan kemampuan penyaringan gangguan di Bab 8: ditambahkan penjelasan praktis tentang identifikasi penelepon, daftar hitam/putih, filter kata kunci SMS, senyapkan panggilan tak dikenal di iOS, keterkaitan mode Jangan Ganggu, dan lainnya, serta saling tertaut dengan Bab 13.
  \item Saran perlindungan tiga lapis: penyaringan sistem + anti-gangguan operator + National Anti-Fraud Center, dengan menekankan daftar putih agar kontak penting tidak terblokir salah.
\end{itemize}

\booksection{Tambahan Konten: Peringatan Cuaca / Peringatan Bencana Alam / Peringatan Gempa}

\begin{itemize}
  \item Bab 13 menambahkan bagian ketujuh: peringatan cuaca (kanal resmi meteorologi + App cuaca bawaan sistem + tingkatan warna), peringatan bencana alam (berbagai kanal meteorologi/tanggap darurat), serta prinsip dan miskonsepsi peringatan gempa (peringatan dini $\neq$ prakiraan).
  \item Tabel pengaktifan peringatan gempa per merek: jalur umum dan catatan soal izin untuk Xiaomi, Huawei/HONOR, keluarga OPPO, vivo, iPhone, Samsung, dan lainnya.
  \item Bab 32 menambahkan "Skenario Nol": urutan menyelamatkan diri setelah peringatan gempa, penanganan hujan ekstrem/taifun, serta keterkaitan dengan SOS satelit / komunikasi tanpa jaringan; panduan P3K dan pengaturan peringatan di Bab 13 saling tertaut.
  \item Keterkaitan subbab cuaca: entri App cuaca pada bagian perjalanan diperkaya dengan notifikasi peringatan serta pengingat "notifikasi harus diaktifkan agar peringatan bisa diterima".
\end{itemize}

\booksection{Tambahan Konten: Kolaborasi Lintas Perangkat dan Alat Produktivitas}

\begin{itemize}
  \item Bab 13 bagian kedelapan: mengendalikan komputer dari ponsel (Sunflower, ToDesk) dan batas merah anti-penipuan; kantor AI (Trae / WorkBuddy / Dumate) dan alur kerja mengendalikan komputer dari ponsel; serta wallpaper dinamis Artificial Desktop dari miHoYo.
  \item Kirim cepat jaringan lokal dan papan klip: LocalSend, Feige Chuanshu, AirDrop; KDE Connect, papan klip universal, serta transfer antara ponsel dan komputer satu merek.
  \item Kata sandi dinamis 2FA: Aegis, Microsoft Authenticator, kode cadangan, dan hal yang perlu diperhatikan saat ganti ponsel.
  \item TickTick dan pemindaian kertas: perincian penggunaan untuk jadwal/tugas; dilengkapi alat pemindai dan teknik memotret dokumen.
\end{itemize}

Semoga kamu makin cakap menggunakan ponsel, menikmati kemudahan dan keseruan yang dibawa teknologi, serta menjadi warga digital yang sadar dan percaya diri.

— Pemandu Kehidupan Digitalmu
